\chapter{Eine gemeinsame Ausführung für Record, Präparation und Mehrzeit}
\label{ch:execution13}
\begin{bild}{Aktualisierung v1.5: weniger zusätzliche Steuerung}
Die hier eingeführte gemeinsame Record-/Präparations-/Mehrzeitausführung lässt sich jetzt bei festem $\Delta>0$ ohne negative Evolutionszeiten anschließen. Endliche Q-Gatezeiten sind unter einem ausdrücklich belegungserhaltenden Gatevertrag mitgerechnet. Wedge-Kopplung, Q, adressiertes An/aus und Reset bleiben zusätzliche Ressourcen; Details in Kapitel~\ref{ch:execution15}.
\end{bild}
\section{Eine Konstruktion, aber ein ausdrücklich zusätzlicher Ressourcenvertrag}
Dieselbe korrigierte Vermittlungsoperation erzeugt jetzt Aufzeichnung, kontrollierte Näherungspräparation und Mehrzeitantwort. Die tatsächlichen Quellprojektoren, Paulioperationen und die Dreierclock wurden benutzt. Ein bereits präpariertes $\Omega$ oder ein separater Präparations-Hamiltonoperator wird nicht benötigt.

Gesetzt bleiben gekoppelte Träger, Wedge-Kopplung und ihre unitäre Fortsetzung, kohärenter Belegungszugriff, Registerinitialisierung und -messung, Kantensteuerung und Feedback. Die Messregel ist Voraussetzung. Die gemeinsame Konstruktion ist daher bedingt, nicht vollständig nativ.

\section{Die vollständige Recordoperation}
Mit $W^\dagger W=P_-$, $WW^\dagger=I_6$ und $WP_+=0$ gilt
\begin{equation}
U_0=\begin{pmatrix}P_+&W^\dagger\\W&0\end{pmatrix},\quad U_0^2=I_{22}.
\end{equation}
Kopiere nur die Vermittlerbelegung, nicht das geordnete Farbwort:
\begin{equation}
Q=I_{16}\otimes I_2\oplus I_6\otimes X.
\end{equation}
Dann ist auf allen 44 Dimensionen
\begin{equation}
(U_0\otimes I)Q(U_0\otimes I)
=(P_+\otimes I+P_-\otimes X)\oplus I_{12}.
\end{equation}
Der Vermittler kehrt zurück. Auf Materie und Pointer bleibt $R=P_+\otimes I+P_-\otimes X$, $R^2=I$, mit
\begin{equation}
R(\psi\otimes\ket0)=P_+\psi\otimes\ket0+P_-\psi\otimes\ket1.
\end{equation}
Beide Krausoperatoren gehören zum Instrument; ihre Effekte summieren sich zu $I$. Ein ungemessener Pointer bleibt kohärent.

\section{Ein gemeinsamer Träger und ein fester Schritt}
Für vier Materieträger und höchstens einen aktiven Vermittler sei
\begin{equation}
\mathcal H_{\mathrm{common}}=(\C^4)^{\otimes4}\oplus
\bigoplus_{e=1}^{6}(\C^6_e\otimes\C^4\otimes\C^4),\quad \dim=832.
\end{equation}
Alle sechs $W_e$ wurden auf dem vollen Materieraum rekonstruiert:
$W_e^\dagger W_e=P^-_e$, $W_eW_e^\dagger=I_{96}$, $W_eP^+_e=0$.
Andere Vermittlersektoren bleiben bei einer adressierten Operation unverändert. Die Kanten haben orthogonale Ausgänge; verschränkte Zuschauer sind eingeschlossen. Dies ist ein vollständiges sequenzielles Protokollmodell, kein Fockraum beliebig vieler aktiver Vermittler.

Die tatsächliche Clock $C$ permutiert $0\to1\to2\to0$ und fixiert 3. Sie erhält alle 240 phasenmarkierten Wurzeln und 60 Projektoren und normalisiert die Quell-Paulis. Mit $e(0)=01$, $e(1)=02$, $e(2)=03$ und $R_{e(3)}=I$ gilt
\begin{equation}
\mathbb U=(C\otimes I)\sum_{k=0}^{3}\ket k\bra k\otimes R_{e(k)}.
\end{equation}
Orthogonale Steuerblöcke machen $\mathbb U$ unitär. Mit Clock und Pointer hat der Raum Dimension 6656. Die Clock-Kanten-Zuordnung und das spätere Austauschen oder Auslesen der Pointer bleiben gesetzte Zugriffe. Eine physikalische Sekunde entsteht nicht aus der Dimensionszahl.

\section{Quellseitiger Reset ohne vorpräpariertes Zielreservoir}
Rechenbasis- und uniformer Projektor gehören zu den 60 tatsächlichen Quellprojektoren. Die Swap-Identität verwendet dieselben 16 Quell-Paulis:
\begin{equation}
S=\frac14\sum_{v=0}^{15}P_v\otimes P_v.
\end{equation}
Für Messausgang $x$ und Ziel $a$ liefert die Bitverschiebung
\begin{equation}
M_x=X_{a\oplus x}\ket x\bra x=\ket a\bra x,\quad
\sum_xM_x^\dagger M_x=I.
\end{equation}
Damit ist $\sum_xM_x\rho M_x^\dagger=\ket a\bra a\,\Tr\rho$. Vier lokale Resets erzeugen $v_0=\ket{0123}$ sogar aus verschränktem Eingang. Messung und Feedback bleiben Ressourcen, keine hergeleitete kosmologische Zustandswahl.

\section{Wiederholung präpariert den Viererzustand kontrolliert}
Akzeptiere in jeder Sternrunde die Recordfolge 111:
\begin{equation}
K_\star=P^-_{03}P^-_{02}P^-_{01},\quad v_N=K_\star^Nv_0,\quad p_N=\|v_N\|^2.
\end{equation}
Alle acht ersten Recordfolgen haben Wahrscheinlichkeit $1/8$. $\Omega$ wird erst zur Zielprüfung definiert. Im 24-dimensionalen Raum der verschiedenen Farbanordnungen gilt exakt
\begin{equation}
\det(xI-K_\star^\dagger K_\star)=
\frac{x^{12}(x-1)(16x-1)^5(16x^2-9x+1)^3}{2^{32}}.
\end{equation}
Dieser Raum trägt die treue reguläre Darstellung von $S_4$. Die polynomialen Identitäten der Gruppenalgebra gelten daher auch auf allen 256 Materiedimensionen. Die übrigen Eigenmultiplizitäten werden dabei nicht unverändert übernommen. Der Einsprojektor ist $P_\Omega$. Weil $K_\star$ und sein Adjungiertes ihn fixieren, folgt mit $\beta=(9+\sqrt{17})/32<5/12$
\begin{align}
\|K_\star^N-P_\Omega\|^2&\le\beta^N,\\
1/24\le p_N&\le1/24+(23/24)\beta^N,\\
F_N=1/(24p_N)&\ge1/(1+23\beta^N).
\end{align}
\begin{center}\begin{tabular}{rll}
\toprule $N$&Erfolgswahrscheinlichkeit&Bedingte Zielfidelität\\\midrule
1&$1/8$&$1/3$\\
2&$51/1024$&$128/153$\\
4&$175341/4194304$&$524288/526023$\\
8&$\approx0{,}04166669813$&$\approx0{,}9999992449$\\
$\infty$&$1/24$&$1$\\\bottomrule
\end{tabular}\end{center}
Exakt sind $p_8=2932033221441/70368744177664$ und
$F_8=8796093022208/8796099664323$. Der konkrete Eingang erreicht bei acht Runden Infidelität unter $10^{-6}$; die allgemeine rationale Schranke garantiert dies bei 20 Runden. Endliches $N$ ist keine exakte Zielpräparation.

Fehlschläge lösen einen erklärten Reset aus. Wegen $p_N\ge1/24$ benötigt man höchstens 24 Versuche im Mittel; $m$ Fehlschläge haben Wahrscheinlichkeit höchstens $(23/24)^m$, bei $m=500$ unter $10^{-9}$. Beendigung ist fast sicher, nicht deterministisch nach einer festen Versuchszahl. Die frühere Überlappung $1/6$ setzt den Zweisingulett-Eingang voraus: Seine Erzeugung aus $v_0$ durch Paarfilter kostet $1/4$, insgesamt wieder $1/24$.

\section{Derselbe Filter liest die Mehrzeitantwort aus}
Nach erfolgreicher Präparation folgen lokaler Quelltick $A$, zweimal $R_{01}$, Rücknahme $A^\dagger$ und dasselbe $N$-Runden-Filter. Ein ideales kostenloses $\Omega$-Messgerät wird nicht eingesetzt. Protokoll R behält denselben kohärenten Pointer; Protokoll F verwendet beim zweiten Schritt einen frischen. Mit $B_\pm=A^\dagger P^\pm_{01}A$ lauten die gemeinsamen unbedingten Erfolge
\begin{equation}
j_N^R=\|K_\star^{2N}v_0\|^2,\qquad
j_N^F=\sum_{\sigma=\pm}\|K_\star^NB_\sigma K_\star^Nv_0\|^2.
\end{equation}
Division durch $p_N$ ergibt den bedingten Schlusserfolg. Die Zweige werden als Dichten, nicht kohärent als Amplituden addiert.

Für die tatsächliche Dreierclock mit Spur eins sind die idealen Recordgewichte $5/8,3/8$. Bedingte Grenzwerte sind $1,17/32$, unbedingte $1/24,17/768$.
\begin{center}\begin{tabular}{rrr}
\toprule $N$&Record behalten&Record frisch\\\midrule
1&0,3984375&0,25\\
2&0,8393698300&0,4434886259\\
4&0,9967024178&0,5274997175\\
8&0,9999992449&0,5312581678\\
12&0,9999999998&0,5312504084\\
$\infty$&1&0,53125\\\bottomrule
\end{tabular}\end{center}
Die Tabelle zeigt bedingte Schlusserfolge; intern wurden exakte Brüche gerechnet. Die Clock erzeugt wiederholte Farben und verlässt den 24-dimensionalen Präparationssektor. Deshalb wurde das Echo auf allen 256 Dimensionen geprüft.

Der andere Quelltick $Z\otimes I$ hat Spur null und liefert $1$ gegen $5/9$. Allgemein gilt $p_+=(16-|\Tr A|^2)/24$. Die ältere ausgeglichene Variante verwendet Spur zwei und liefert $1$ gegen $1/2$. Unterschiedliche Tick-Settings werden nicht vermischt.

\section{Was für eine native Implementierung noch fehlt}
Das neue Register ist nicht Clifford:
\begin{equation}
R(I_{16}\otimes Z)R^\dagger=S\otimes Z,\quad \Tr S=4.
\end{equation}
Ein vierqubitiger Pauli hat Spur null oder $\pm16$. Noch konkreter erzeugt der Quellinput $\ket0\otimes(\sum_b\ket b/2)\otimes\ket0$ genau 13 nichtverschwindende Rechenbasisamplituden. Das ist kein reiner Stabilizerzustand, dessen Trägergröße eine Zweierpotenz ist. Nur Cliffordoperationen, Stabilizer-Hilfszustände und Pauli-Messungen reichen auch mit Feedback nicht aus. Andere Kodierungen oder Kopplungen wären zusätzliche Ressourcen, kein pauschal ausgeschlossener TFPT-Weg.

Die zentrale Phase $iI_4$ wirkt sowohl auf zwei Eingängen als auch auf dem Wedge-Vermittler als minus Identität. Sie liefert nicht den relativen Belegungspuls $D_{\mathrm{occ}}=I_{16}\oplus(-I_6)$. Dieser erfüllt $U_0D_{\mathrm{occ}}U_0=S\oplus I_6$, muss aber als eigener Zugriff begründet sein.

\section{Ungelesener Austausch ist keine Kühlung}
Alle $P^\pm_e$ kommutieren mit $P_\Omega$; der ungelesene Kanal erhält das Zielgewicht. Seine einzelnen Dephasierungen sind orthogonale Hilbert--Schmidt-Projektionen. Ein Fixpunkt ihres Zyklus kann bei keinem Schritt Norm verlieren und liegt deshalb im Schnitt der Fixräume: in der Kommutante der $S_4$-Wirkung. Deren Dimension ist
\begin{equation}
\dim\mathrm{Fix}=\frac1{24}\sum_{\pi\in S_4}4^{2c(\pi)}=3876.
\end{equation}
Die spurlose Dimension 3875 identifiziert keinen gleichnamigen $E_8$-Modul. Die frühere Zahl 23076 betrifft die größere Kommutante eines einzelnen Hamiltonoperators. Reset und Feedback umgehen die Gewichtserhaltung durch zusätzliche Operationen; sie liefern keinen intrinsischen kosmologischen Attraktor.

\section{Abnahmetests und T1--T8-Grenze}
Ein nativer Adapter muss $W$, $W^\dagger$, Kern und Belegung gemeinsam liefern, beim Half-Charge-Feld zusätzlich mit Domänen- und Energiekontrolle. Er muss den 13-Amplituden-Eingang, beide Recordzweige und die Acht-Runden-Echos reproduzieren, einschließlich Hilfszuständen und Fehlerausgängen. Erst anschließend ist eine skalierende Zellfamilie sinnvoll auf Geometrie, Chiralität und Spin 2 zu prüfen.

T1 betrifft weiter die primitive Auswahl; T2 die native analytische Naht; T3 den gemeinsamen 3+1D-Ursprung; T4 das chirale Maß; T5 den wechselwirkenden Grenzprozess; T6 vollständige Parameter und Transfer; T7 dynamischen Spin 2; T8 das physikalische Quellfunktional. Alle vollständigen Tore bleiben offen. Die 231 eigenen exakten Bedingungen und 1073 getrennten Quell-Präfixprüfungen wurden für Ausgabe v1.3 erneut ausgeführt; Normal- und optimierter Lauf stimmen bytegenau überein, fünf Negativmutanten wurden erkannt. Dies ist die Provenienz jener Ausgabe, keine Behauptung eines weiteren Laufs dieser Suite in v1.4.
